% Source: https://artofproblemsolving.com/community/c3h2669400_111th_post_solving_problems
% Stable identity: aops-topic-2669400-q44
Write $T_k=\dfrac{k(k+1)}{2}$. For $s=1$ take $x_1=T_1=1$.

\emph{Small cases.} $\dfrac13+\dfrac13+\dfrac13=1$ with $3=T_2$, and since $\dfrac13=\dfrac16+\dfrac16$ with $6=T_3$, also
\[\dfrac13+\dfrac13+\dfrac16+\dfrac16=1,\qquad\dfrac13+\dfrac16+\dfrac16+\dfrac16+\dfrac16=1.\]
This settles $s=3,4,5$.

\emph{Adding three terms.} For every $k$,
\[\dfrac1{T_{2k}}+\dfrac1{T_{2k+1}}=\dfrac{2}{2k+1}\left(\dfrac1{2k}+\dfrac1{2k+2}\right)=\dfrac{2}{2k+1}\cdot\dfrac{4k+2}{2k(2k+2)}=\dfrac{1}{k(k+1)}=\dfrac1{2T_k}.\]
So one term $\dfrac1{T_k}$ can be replaced by the four terms $\dfrac1{T_{2k}}+\dfrac1{T_{2k+1}}+\dfrac1{T_{2k}}+\dfrac1{T_{2k+1}}$, increasing $s$ by $3$. Starting from $s=3,4,5$, we reach every $s\ge3$. $\blacksquare$
